\originalsection{第9次习题课}
\sourceinfo{MIT 18.04, Spring 2018, Recitation 9; 题面 \nolinkurl{mit18_04s18_recit9-handout.pdf}, PDF第1页; 解答 \nolinkurl{mit18_04s18_recit9-solutions.pdf}, PDF第1--2页. 课程作者 Jeremy Orloff, 页内署名 Vishesh Jain, 2018, CC BY-NC-SA 4.0}
\begin{exercise}
\textbf{原题 Rec9.1。}\label{ass:rec9-1}
计算 $I_1=\int_{-\infty}^\infty(1+x^4)^{-1}\dd x$.
\end{exercise}
\begin{solution}
编者补充 (AI): 官方解答PDF第1页仅指向讲义例9.4. 对 $f(z)=1/(1+z^4)$ 用上半圆闭轮廓, $R>1$. 内部两极点为 $a=\ee^{\ii\pi/4}$ 与 $b=\ee^{3\ii\pi/4}$, 均简单, 留数为 $1/(4a^3)$ 与 $1/(4b^3)$. 它们之和为 $(\ee^{-3\ii\pi/4}+\ee^{-9\ii\pi/4})/4=-\ii/(2\sqrt2)$. 留数定理给闭路积分为 $2\pi\ii(-\ii/(2\sqrt2))=\pi/\sqrt2$. 圆弧上 $|1+z^4|\ge R^4-1$, 圆弧积分绝对值不超过 $\pi R/(R^4-1)\to0$. 实积分的尾部为 $O(x^{-4})$, 两端均收敛, 因而 $I_1=\pi/\sqrt2$.
\end{solution}
\begin{exercise}
\textbf{原题 Rec9.2。}\label{ass:rec9-2}
计算 $I_2=\int_0^{2\pi}(2-\sin\theta)^{-1}\dd\theta$.
\end{exercise}
\begin{solution}
据官方解答翻译补全 (解答PDF第1页). 令 $z=\ee^{\ii\theta}$, 则 $\dd\theta=\dd z/(\ii z)$, $\sin\theta=(z-z^{-1})/(2\ii)$. 故
\[
I_2=\int_{|z|=1}\frac{-2}{z^2-4\ii z-1}\dd z.
\]
分母根为 $a=\ii(2-\sqrt3)$ 与 $b=\ii(2+\sqrt3)$, 仅 $a$ 在单位圆内. 在 $a$ 处留数为 $-2/(a-b)=1/(\ii\sqrt3)$, 因而 $I_2=2\pi\ii/(\ii\sqrt3)=2\pi/\sqrt3$. 原实分母至少为1, 所以不存在路径上的奇点.
\end{solution}
\begin{exercise}
\textbf{原题 Rec9.3。}\label{ass:rec9-3}
计算 $I_3=\int_0^{2\pi}(\sin\theta)^{2n}\dd\theta$. 原题没有声明 $n$ 的范围; 以下取 $n\in\Z_{\ge0}$.
\end{exercise}
\begin{solution}
据官方解答翻译补全 (解答PDF第1--2页); $n$ 的范围为编者补充 (AI). 同上代换得
\[
I_3=\int_{|z|=1}\frac{(z^2-1)^{2n}}{2^{2n}\ii^{2n+1}z^{2n+1}}\dd z.
\]
只有0是极点, 留数为分子中 $z^{2n}$ 的系数除以 $2^{2n}\ii^{2n+1}$. 二项式展开给该系数 $\binom{2n}{n}(-1)^n$, 又 $\ii^{2n}=(-1)^n$, 所以留数为 $\binom{2n}{n}/(2^{2n}\ii)$. 因此 $I_3=2\pi\binom{2n}{n}/4^n=2\pi(2n)!/(2^{2n}(n!)^2)$. 当 $n=0$ 时仍给出 $2\pi$, 与被积函数恒为1相符.
\end{solution}
\begin{exercise}
\textbf{原题 Rec9.4。}\label{ass:rec9-4}
计算 $I_4=\int_1^\infty\dd x/(x\sqrt{x^2-1})$.
\end{exercise}
\begin{solution}
编者补充 (AI): 官方解答PDF第2页仅指向讲义例9.8. 对 $x>1$ 的根号取正实值. 端点1附近被积函数为 $O((x-1)^{-1/2})$, 无穷远为 $O(x^{-2})$, 所以广义积分收敛. 令 $x=\sec t$, $0<t<\pi/2$, 则 $\sqrt{x^2-1}=\tan t$ 且 $\dd x=\sec t\tan t\dd t$, 故 $I_4=\int_0^{\pi/2}\dd t=\pi/2$. 这个实变量代换同时明确了根号分支和两个端点的取极限方向.
\end{solution}
